\section{Benchmark Construction and Evidence}
\label{sec:benchmark}
\subsection{Source provenance and functional context}
\bench{} begins from traceable computations rather than arbitrary code padded
with unrelated lines. Each mother case records its source revision, adapted
mathematical definition, public contract, classical implementation, and private
checking evidence. Source tasks and the newly written application contexts are
identified separately. An adaptation from a quantum programming task is not
presented as an original production application supplied by its upstream authors.

Table~\ref{tab:cases} lists development cases and obligations. Shared mathematical
structure and upstream tasks prevent treating every directory or narrative as an
independent problem family.

\begin{table}[tb]
\caption{Development case inventory. IDs identify cases, not independent sample
counts. Scope-control labels remain provisional. Corpus measurements are unfilled.}
\label{tab:cases}
\begin{tabularx}{\linewidth}{lYYY}
\toprule
Case & Source lineage & Classical role & Contract focus \\
\midrule
lit-001 & C2$|$Q$\rangle$ & Weighted partitioning & Weights, canonical selection, change report \\
lit-002 & C2$|$Q$\rangle$ & Vertex-cover work list & Cardinality, tuple ties, ownership \\
lit-003 & C2$|$Q$\rangle$ & Agenda coloring & Preview versus completion, first solution \\
lit-004 & C2$|$Q$\rangle$ & Compatible records & Filtering, recency, maximum subset \\
lit-005 & QuanBench / QuanBench+ & Feature configuration & Locked values, validation order \\
lit-006 & QuanBench & Transfer allocation & Capacity, windows, offsets, ties \\
lit-007 & SupermarQ & Signed preferences & Sign conventions, scores, movement \\
lit-008 & Qiskit HumanEval & Ordered XOR receipts & Full output, order, prefix checksum \\
lit-009 & HPL & Dense balance solving & Pivot selection, exceptions, residuals \\
lit-010 & HPCG & Sparse iterative solving & Initial state, stopping, full trajectory \\
\bottomrule
\end{tabularx}
\end{table}

C2$|$Q$\rangle$, QuanBench, QuanBench+, Qiskit HumanEval, and SupermarQ supply
task or objective provenance~\cite{c2q,quanbench,quanplus,qhe,supermarq}.
HPL and HPCG supply numerical algorithm references~\cite{hpl,hpcg}.
These uses do not reproduce the upstream systems' complete workloads or official
performance measurements. Circuit resources and representation levels can be
studied using ideas from MQT Bench~\cite{mqt}, but resource extraction does not
establish the semantic correctness of the enclosing application.

\subsection{Controlled input views}
For each mother case, view A contains the computational core; view B contains
the full classical program plus a source-location cue; view C contains the same
full program without that cue. The B/C pair is byte-identical after removal of
the designated cue. This makes B/C the clean contrast for location assistance.
A/B changes both context and the obligations visible to the model, so differences
cannot be attributed solely to search difficulty. An analysis of core obligations
shared across A and B is reported separately from full-contract completion.

Context must affect the public contract, as assessed through dependency and
mutation checks. Program size, branches, obligations, and dependency distance are
descriptive covariates; difficulty is measured rather than assumed.

Public input packages contain only allowlisted source and task material. Private
reference proposals, tests, previous responses, and diagnostic conclusions are
excluded. Source notices preserve attribution and licensing, so the benchmark
does not claim perfect source anonymity. Any redaction or renaming experiment
must retain legally required notices and use a separately versioned input.

\subsection{Reference construction and uncertainty}
The reference process begins with a behavioral inventory: input domain, returned
values, ordering, exception behavior, relevant state, and downstream uses. Reviewers
then identify possible regions and admissible formulations. They record whether
a concrete mapping is supported, contradicted, or unresolved and state the scope
of each judgment. Multiple valid regions and different formulation families may
be accepted when their obligations are met.

The current materials contain AI-assisted proposals and coordinator review;
these are not independent gold annotations. For a confirmatory release, separate
reviewers should inspect the original program and contract before seeing model
answers. The proposed process records identity or role, expertise, independence,
initial judgments, disagreements, and adjudication. Difficult cases remain
uncertain when evidence is insufficient. The final inclusion policy is a human
scientific decision, recorded before confirmatory responses are collected.

Correctness and completion occupy different axes. A formula explicitly supplied
but refuted is \emph{provided and contradicted}. A response that honestly defers
the tie-breaking construction is \emph{unresolved and incomplete}. Neither should
be converted into the other by a keyword score. Obligations cover core
correspondence, encoding, exact selection or certification, decoding, and context
preservation. Each record binds the original response passage, its interpretation,
the applicable obligation, evidence, and remaining uncertainty.

\subsection{Local checkers and their scope}
Trusted checkers compare a supported claim against a definition of the classical
task. For a QUBO, a bounded exhaustive checker enumerates assignments and checks
\emph{every} globally optimal decoded state; accepting one convenient ground
state would miss invalid tied optima. For a search predicate, it compares marked
states with the task relation and separately handles selection order, illegal
encodings, and no-solution behavior. Numeric checks preserve the declared
arithmetic model and distinguish permitted inputs from reachable intermediate
states.

The checker suite includes known-correct constructions, equivalent encodings,
and deliberate semantic errors. Examples include reversed optimization direction,
missing constraints, wrong tie order, discarded weights, and omitted iteration
trajectories. These controls validate an evaluator; they are never counted as
model errors or model repairs. A checker can also fail: equality of the zero-energy
set alone, for example, does not establish that feasible states minimize the
energy. The validation record must include extremum checks and retain the failure
that motivated a correction.

Each verdict names its representation, tested domain, arithmetic, and executable
evidence. Unknown representations and infrastructure errors stay distinct from
semantic refutations. In the existing workflow, some natural-language claims
require a reviewer's structured transcription. This is not an automatic judge of
arbitrary prose. The experimental design measures that transcription burden and
reports what fraction of claims can actually be checked.

\figplaceholder{evidence}{2.1cm}{Four layers: public source/contract; private
reference and admissible alternatives; model claim with exact passage; scoped
check and unresolved obligations. Draw a boundary between development feedback
and evaluator-reserved final tests.}
{Provenance and evaluation separation. A verdict is attached to a specific claim
and domain, rather than promoted to whole-program correctness.}
